%% =================================================================
%% Kingma, Ba.
%% "Adam: A Method for Stochastic Optimization."
%% Published as a conference paper at ICLR 2015.
%% arXiv:1412.6980v9 [cs.LG] 30 Jan 2017.
%%
%% Reconstruction of page 1 (printed p. 1) as study notes.
%% Generated by: reproduce-all-pages-basic-tex via
%%               reproduce-multiple-pages-basic-tex [1, 15]
%%
%% Structure: title page — title, authors + affiliations, abstract,
%% Section 1 INTRODUCTION (first portion). No equations, no figures.
%% Note: the rotated arXiv:1412.6980v9 marker on the left margin was
%% captured by pdftotext at its physical y-position inside the abstract
%% region — it is excluded from the body since it is marginalia, not
%% abstract content.
%% =================================================================

\documentclass[11pt]{article}

\usepackage[margin=1in]{geometry}
\usepackage{amsmath, amssymb}
\usepackage{mathrsfs}
\usepackage{bm}
\usepackage{xcolor}
\usepackage{fancyhdr}

\pagestyle{fancy}
\fancyhf{}
\lhead{\small Published as a conference paper at ICLR 2015}
\rhead{}
\cfoot{\small 1 $\to$ \arabic{page}}
\renewcommand{\headrulewidth}{0.4pt}

\setcounter{page}{1}

\begin{document}

% ---------- Title ----------
\begin{center}
{\Large\scshape Adam: A Method for Stochastic Optimization}\\[1.5em]
\begin{tabular}{cc}
\textbf{Diederik P. Kingma\textsuperscript{*}} & \textbf{Jimmy Lei Ba\textsuperscript{*}} \\
University of Amsterdam, OpenAI & University of Toronto \\
\texttt{dpkingma@openai.com} & \texttt{jimmy@psi.utoronto.ca}
\end{tabular}
\end{center}

\vspace{1em}

% ---------- Abstract ----------
\begin{center}
{\scshape Abstract}
\end{center}
\begin{quote}
We introduce \emph{Adam}, an algorithm for first-order gradient-based
optimization of stochastic objective functions, based on adaptive
estimates of lower-order moments. The method is straightforward to
implement, is computationally efficient, has little memory
requirements, is invariant to diagonal rescaling of the gradients,
and is well suited for problems that are large in terms of data
and/or parameters. The method is also appropriate for non-stationary
objectives and problems with very noisy and/or sparse gradients. The
hyper-parameters have intuitive interpretations and typically require
little tuning. Some connections to related algorithms, on which
\emph{Adam} was inspired, are discussed. We also analyze the
theoretical convergence properties of the algorithm and provide a
regret bound on the convergence rate that is comparable to the best
known results under the online convex optimization framework.
Empirical results demonstrate that Adam works well in practice and
compares favorably to other stochastic optimization methods. Finally,
we discuss AdaMax, a variant of Adam based on the infinity norm.
\end{quote}

\vspace{1em}

% ---------- Section 1: Introduction ----------
\section*{1 \textsc{Introduction}}

Stochastic gradient-based optimization is of core practical importance
in many fields of science and engineering. Many problems in these
fields can be cast as the optimization of some scalar parameterized
objective function requiring maximization or minimization with respect
to its parameters. If the function is differentiable w.r.t. its
parameters, gradient descent is a relatively efficient optimization
method, since the computation of first-order partial derivatives
w.r.t. all the parameters is of the same computational complexity as
just evaluating the function. Often, objective functions are
stochastic. For example, many objective functions are composed of a
sum of subfunctions evaluated at different subsamples of data; in this
case optimization can be made more efficient by taking gradient steps
w.r.t. individual subfunctions, i.e. stochastic gradient descent (SGD)
or ascent. SGD proved itself as an efficient and effective
optimization method that was central in many machine learning success
stories, such as recent advances in deep learning (Deng et al., 2013;
Krizhevsky et al., 2012; Hinton \& Salakhutdinov, 2006; Hinton et al.,
2012a; Graves et al., 2013). Objectives may also have other sources
of noise than data subsampling, such as dropout (Hinton et al.,
2012b) regularization. For all such noisy objectives, efficient
stochastic optimization techniques are required. The focus of this
paper is on the optimization of stochastic objectives with
high-dimensional parameters spaces. In these cases, higher-order
optimization methods are ill-suited, and discussion in this paper
will be restricted to first-order methods.

We propose \emph{Adam}, a method for efficient stochastic optimization
that only requires first-order gradients with little memory
requirement. The method computes individual adaptive learning rates
for different parameters from estimates of first and second moments
of the gradients; the name \emph{Adam} is derived from adaptive
moment estimation. Our method is designed to combine the advantages
of two recently popular methods: AdaGrad (Duchi et al., 2011), which
works well with sparse gradients, and RMSProp (Tieleman \& Hinton,
2012), which works well in on-line and non-stationary settings;
important connections to these and other stochastic optimization
methods are clarified in section~5. Some of \emph{Adam}'s advantages
are that the magnitudes of parameter updates are invariant to
rescaling of the gradient, its stepsizes are approximately bounded by
the stepsize hyperparameter, it does not require a stationary
objective, it works with sparse gradients, and it naturally performs
a form of step size annealing.

\vspace{2em}

\noindent\rule{0.3\textwidth}{0.4pt}\\[0.3em]
{\footnotesize\textsuperscript{*}Equal contribution. Author ordering
determined by coin flip over a Google Hangout.}

% [page ends here — continues on page 2]

\end{document}
